\chapter{Conclusiones}

Esta investigación ha estimado el impacto del cambio climático --canalizado a través del estrés térmico-- en el mercado internacional del café, utilizando un modelo estructural de competencia secuencial (Stackelberg) que distingue entre países exportadores líderes y seguidores. Los resultados revelan cómo las alteraciones climáticas reconfiguran los equilibrios estratégicos y la distribución de rentas en un oligopolio global.

Entre los hallazgos principales, destaca la interpretación del cambio climático como un determinante del poder de mercado en el caso del café. El estrés térmico actúa como un choque de costos heterogéneo que amplifica las asimetrías preexistentes entre líderes y seguidores. Su impacto sobre los costos marginales es diferente entre grupos, e incluso diferente entre países dentro del mismo grupo. Esto sugiere que el cambio climático es un factor que puede exacerbar las desigualdades estructurales en mercados de \textit{soft commodities}.

Adicionalmente, la baja elasticidad-precio de la demanda del café juega un rol crucial en la transmisión de estos efectos. La naturaleza inelástica de la demanda implica que los aumentos de costos inducidos por el cambio climático se trasladan de manera significativa a los precios finales, sin que esto genere una contracción proporcional en el consumo. Este mecanismo tiene dos implicancias distributivas relevantes. Por un lado, permite que los productores con menores incrementos de costos —típicamente aquellos con ventajas geográficas o tecnológicas— capturen parte de las rentas generadas por el alza de precios, mitigando parcialmente sus pérdidas directas e incluso obteniendo ganancias estratégicas. Por otro lado, transfiere parte del costo del cambio climático a los consumidores finales, quienes enfrentan precios más altos sin reducir sustancialmente su demanda. En este contexto, la inelasticidad de la demanda no solo amplifica la volatilidad de precios ante choques de oferta, sino que también redistribuye los efectos del cambio climático entre productores, favoreciendo a algunos, mientras que los consumidores absorben una porción creciente del costo agregado.

En segundo lugar, el escenario factual generó importantes pérdidas en el bienestar social agregado, con precios más altos y un menor volumen comercializado. Sin embargo, esta pérdida agregada oculta redistribuciones contraintuitivas. Entre estas, destaca el caso de Vietnam --un país especializado en Robusta, que es la variedad del grano de café más tolerante al calor--, el que vio incrementada su ventaja competitiva en un escenario con calentamiento, a pesar de su pérdida de eficiencia productiva. Esto plantea un dilema importante en cuanto al diseño de políticas globales de mitigación climática, pues podrían encontrar resistencia de aquellos países que se han visto beneficiados por el aumento del estrés térmico en el tiempo.

Además, el caso de Vietnam muestra que, aunque el calor también incrementa sus costos marginales en el factual, su posición relativa mejora porque perjudica más a sus competidores directos (Brasil y Colombia). En el equilibrio de Stackelberg, esto se traduce en una ganancia de cuota de mercado y excedente para Vietnam. El efecto final sobre un país depende tanto de la heterogeneidad en la exposición climática como de su posición dentro de la jerarquía del mercado.

En cuanto a la contribución a la literatura, ésta se da por tres componentes principales. A la economía ambiental del cambio climático, aporta una metodología estructural que descompone los impactos agregados en efectos de eficiencia y estratégicos, demostrando que las estimaciones de modelos competitivos o reducidos podrían no capturar dinámicas distributivas complejas. A la organización industrial de los mercados agrícolas, provee evidencia de que los choques climáticos constituyen un determinante significativo de la evolución del poder de mercado y las estructuras de competencia en el largo plazo. Por último, a la literatura específica relativa al mercado internacional del café, se contribuye con una aproximación para estimar el impacto del estrés térmico sobre los costos marginales y el equilibrio estratégico o las condiciones competitivas de cada agente.

Ahora bien, las conclusiones principales están sujetas a limitaciones relevantes que abren caminos para la investigación futura. Primero, el supuesto de un bien homogéneo, aunque útil y con fundamento histórico, impide analizar la adaptación mediante un cambio en la variedad del café (de Arábica a Robusta). En ese sentido, un modelo que incorpore ambas variedades como sustitutos imperfectos sería un avance natural. Segundo, la estructura de competencia se asume fija en el tiempo; esto obedece a la evolución de la posición competitiva de los líderes en el tiempo y también a fundamentos históricos, aunque modelar endógenamente la transición de un país seguidor a líder (como el caso de Vietnam) en respuesta a choques climáticos y tecnológicos sería un desafío relevante. Tercero, el análisis captura el impacto del estrés térmico bajo el supuesto de que la distribución geoespacial de los cultivos se mantiene constante en el tiempo, y que es igual a lo observado en el año 2020, es decir, con posterioridad al período analizado. Esto implica que los resultados se interpreten como condicionales a la adaptación acumulada en el período anterior a 2020, y sin posibilidad para los agentes de ajustar endógenamente su estructura productiva en respuesta al estrés térmico.

Con todo, los resultados tienen implicancias concretas para el diseño de políticas públicas. Por un lado, es importante que los programas de asistencia técnica y financiera para la adaptación climática sean diferenciados y estratégicos. Así, priorizar a los países seguidores más vulnerables puede no solo ser una cuestión de equidad, sino también de eficiencia para contrarrestar la concentración indebida del mercado.

Por otro lado, desde una perspectiva de estrategia de desarrollo nacional, los países exportadores deben evaluar su ventaja comparativa climática futura. Para algunos, la estrategia óptima puede ser invertir en la transición a variedades resilientes (Robusta); para otros, como los productores de Arábica de alta calidad, la estrategia puede radicar en invertir en nuevas tecnologías para mitigar el impacto negativo de la exposición al calor.

En síntesis, este trabajo demuestra que el impacto económico del cambio climático en el mercado internacional del café trasciende el cálculo de pérdidas promedio de rendimiento. Al alterar los costos relativos y las posiciones estratégicas dentro de un modelo oligopólico, el cambio climático redistribuye el poder de mercado y las rentas económicas, creando nuevos retos para eficiencia, la equidad y la gobernanza del sector cafetalero.